\section{Results}
\label{sec:results}

\begin{table}[t]
\centering
\small
%% GENERATED by pipeline/latex.py -- do not edit.
%% Regenerate with: python3 pipeline/latex.py
\begin{tabular}{lrr}
\toprule
Injection syntax & CVEs & with PoC \\
\midrule
HTML / DOM (XSS)    &  50{,}884 & 27{,}418 \\
SQL                 &  22{,}225 & 15{,}011 \\
Shell command       &   9{,}672 & 5{,}305 \\
Code / eval         &   6{,}666 & 3{,}972 \\
CRLF / header       &       656 &     331 \\
Argument            &       430 &     216 \\
\bottomrule
\end{tabular}

\caption{Largest injection syntax types, of \dataInjectionCVEs injection CVEs, with
the count carrying at least one linked proof-of-concept.}
\label{tab:syntax}
\end{table}

\paragraph{Syntax distribution.}
\Cref{tab:syntax} shows that injection is overwhelmingly concentrated in two
syntaxes (HTML/DOM and SQL together account for four in five injection CVEs)
with a long tail of shell, code-eval, header, template, LDAP, XML, and
formula contexts.
This concentration matters directly for \cref{sec:prevent}: the two dominant
syntaxes are exactly the ones a matchertext host can delimit with a matcher pair.
Exploitation signals are sparse but consistent with severity: \dataInjectionKev injection CVEs
appear in the CISA Known-Exploited catalog, and command injection, though
only a tenth of the corpus, carries a disproportionate share of them.

\paragraph{Proof-of-concept coverage.}
\textbf{\dataInjectionWithPoc} injection CVEs (\dataInjectionWithPocShare) carry at least one linked PoC.
The PoC databases overlap, so unifying them raises payload coverage well
above any single source. Of these, \textbf{\dataPayloadBackedCVEs} yield an extractable payload; the remainder
are links only or exploits from which no literal attack string can be recovered, such as modules
that construct their payload programmatically. This linkage is what makes the prevention analysis
of \cref{sec:prevent} possible, since it supplies the real attack strings.

\paragraph{Sub-classification.}
Beyond the two primary axes, each injection CVE is labelled on orthogonal
facets, each from its strongest available signal, and stored in long form so a
CVE can carry one label per dimension at once.
The CVSS vector yields reachability for free: \textbf{\dataSubPrivilegePreAuth} injection CVEs
are exploitable pre-authentication, \dataSubPrivilegeLowPriv require low privilege, and \dataSubPrivilegeHighPriv
high privilege.
Finer \emph{technique} labels come from PoC payloads and descriptions where
present: SQL injections resolve into blind-boolean (\dataSubTechniqueBlindBoolean), union-based (\dataSubTechniqueUnionBased),
blind-time (\dataSubTechniqueBlindTime), and authentication-bypass (\dataSubTechniqueAuthBypass) attacks, while XSS resolves
into stored (\dataSubTechniqueStored), reflected (\dataSubTechniqueReflected), and DOM-based (\dataSubTechniqueDomBased).

\paragraph{Syntactic groups.}
For the \textbf{\dataPayloadBackedCVEs} injection CVEs with an extractable PoC payload, we
reduce each payload to a \emph{syntactic skeleton}: string literals, numbers,
and identifiers are abstracted to placeholders (\verb|<q>|, \verb|<n>|,
\verb|<id>|) while the structural grammar (matchers, operators, tags, and a
curated set of injection keywords) is preserved.
Grouping by identical skeleton yields \textbf{\dataDistinctSkeletons} distinct groups
(\dataSingletonGroups of them singletons), so specific values never split a group but a
different breakout structure always does.
The largest groups are recognizable attack families: the XSS skeletons
\verb|<TAG>ALERT(DOCUMENT.<id>)</TAG>| (\dataGroupAlertDocId CVEs) and its
\verb|DOCUMENT.COOKIE| variant (\dataGroupAlertCookie), and the code-eval fragment
\verb|<? <id>| (\dataGroupPhpOpen).
The tail is long: the singletons are \dataSingletonGroupShare of groups but only \dataSingletonCVEShare of CVEs, and
they collapse to fifteen distinct verdicts, since the skeleton preserves tags and
operators that the containment test does not depend on.
Because the skeleton preserves exactly the matcher structure while discarding
values, it is the natural unit on which to test matchertext prevention next: the
verdict depends only on that structure, so it is shared by every CVE in a group.
